% Abschnitte formatieren (150-200 Wörter pro Abschnitt)
% Introduction zu nah an der aus der Abschlussarbeit --> Fokus mehr auf Evaluation legen, weniger auf technische Entwicklung
% Beitrag des Papers eher Evaluation als technische Entwicklung
% Ausarbeitung eines Frameworks unpassend --> eher Evaluations Setup basierend auf einer Case Study: Guideline für das Evaluieren von RAG-Systemen + Was sind Learnings für RAG-Systeme in Bildungskontexten
% Forschungsfrage klarer hervorheben (auf Evaluation bezogen) --> contributions daraus ableiten

\section{Introduction}
The digital transformation is fundamentally changing the design of learning environments and increasingly placing intelligent information systems at the center of modern educational platforms. Such systems are understood as socio-technical artifacts that integrate technological components, data structures, and usage contexts to make knowledge accessible and usable \autocite{Klesel2025}. Particularly in the context of digital learning platforms, Artificial Intelligence (AI) not only support the provision of information but can also actively structure and influence learning processes \autocite{ZawackiRichter2019}. Conversational systems relying on Large Language Models (LLMs) have taken on a central role in this regard, as they enable natural language interaction and thus open up new forms of learning support \autocite{Oates2025ChatGPTClassroom}. At the same time, generative models have fundamental limitations, particularly regarding the reliability of the content they generate. Hallucinations pose a critical risk, especially in educational contexts where accurate and verifiable information is essential \autocite{Feuerriegel2024}. Retrieval-Augmented Generation (RAG) addresses this problem by integrating external knowledge sources into the generation process, thereby establishing a link between generative AI and organization-specific knowledge \autocite{Lewis2020_RAG,Klesel2025}.

This paper is situated within the context of the KI-Campus platform, a large-scale, publicly accessible learning environment for AI with over 130,000 users as of early 2026. A RAG-based chatbot is already deployed on the platform to support users in navigating course offerings and answering domain-specific questions. Despite its productive use, the current system exhibits notable shortcomings, particularly with regard to response quality and the lack of dedicated learning support. Furhtermore, there is currently no generally accepted standard for evaluating and comparing RAG-based chatbots \autocite{Yu2024RAGEvalSurvey,Swacha2025}. Against this background, this paper addresses the following research question: How can RAG-based educational chatbots be evaluated and compared using a multidimensional evaluation approach that captures both technical performance and user-centered quality dimensions. To address this question, the study conducts a multidimensional comparative evaluation of two RAG system designs deployed on the KI-Campus platform. The evaluation is based on a case-study-driven setup that combines automated metrics, expert evaluation, user blind assessment, and cost analysis using real interaction data. The goal is to derive practical guidelines for the evaluation of RAG systems and to identify key design trade-offs that arise in real-world educational applications. The main contribution of this work is an application-oriented, multidimensional evaluation setup for RAG-based educational chatbots, demonstrating the necessity of combining diverse assessment perspectives. As an additional contribution, the study offers empirical insights into how different architectural choices affect not only technical performance and costs, but crucially, the pedagogical quality and didactic value in real-world learning scenarios. Based on the empirical results, the paper derives lessons learned and recommendations for the evaluation and design of RAG systems in educational contexts.
